
O paradigma de geração de texto condicionado consolidou-se com o modelo T5 (\textit{Text-to-Text Transfer Transformer})~\cite{raffel_exploring_2020}, que unificou tarefas diversas de PLN — sumarização, tradução e geração de respostas — em um único framework de entrada e saída textual, viabilizando o ajuste fino de modelos pré-treinados para tarefas de geração como a proposta neste trabalho. Brown et al.~\cite{brown_language_2020}, com o GPT-3, demonstraram que LLMs de grande escala podem realizar tarefas de geração via \textit{in-context learning} sem ajuste fino, mas também evidenciaram que, para domínios específicos, o zero-shot produz resultados inferiores aos obtidos com especialização supervisionada — observação corroborada pelo desempenho do \textit{baseline} neste estudo.

A arquitetura \textit{Transformer}~\cite{vaswani_attention_2017} e modelos derivados como o BERT~\cite{devlin_bert_2019} estabeleceram que representações contextualizadas pré-treinadas podem ser ajustadas para tarefas específicas com alta eficiência. O Sentence-BERT (S-BERT)~\cite{reimers_sentence-bert_2019} estende o BERT com uma rede siamesa para gerar \textit{embeddings} de sentenças comparáveis por similaridade semântica, tornando-se ferramenta padrão para avaliação de qualidade em geração de texto — papel que ocupa também neste trabalho.

Com o crescimento dos LLMs, o treinamento completo tornou-se computacionalmente inviável na maioria dos cenários práticos. Ding et al.~\cite{ding_parameter-efficient_2023} apresentam uma revisão sistemática dos principais métodos PEFT — adaptadores, \textit{prefix-tuning} e LoRA — demonstrando que o ajuste eficiente alcança desempenho comparável ao treinamento completo em tarefas de classificação e geração com uma fração dos parâmetros treináveis. O LoRA~\cite{hu_lora_2022} e sua extensão com quantização Q-LoRA~\cite{dettmers_qlora_2023} são os representantes mais utilizados desse paradigma, sendo ambos avaliados neste trabalho.

No contexto nacional, Souza et al.~\cite{souza_bertimbau_2020} demonstraram que o pré-treinamento adicional do BERT em corpora de português brasileiro (BERTimbau) gera ganhos significativos em tarefas como reconhecimento de entidades nomeadas e similaridade textual em relação a modelos multilíngues, evidenciando que a especialização de modelos pré-treinados em um domínio ou idioma é necessária para maximizar o desempenho. Esse princípio motiva diretamente a avaliação de estratégias de ajuste fino conduzida neste trabalho, agora aplicado ao domínio de textos científicos em inglês do arXiv.

Os modelos Qwen2~\cite{yang_qwen2_2024} e Qwen2.5~\cite{team_qwen25_2024} integram a nova geração de LLMs abertos com arquitetura eficiente, apresentando resultados competitivos em benchmarks de compreensão e geração de texto mesmo nas versões de menor porte. Estudos recentes sobre os efeitos da quantização em LLMs~\cite{shi_systematic_2025} indicam que, embora reduza custos computacionais, a quantização pode introduzir degradação de qualidade variável dependendo da tarefa e da arquitetura — aspecto diretamente relevante para a interpretação dos resultados obtidos com o Q-LoRA neste estudo.
